% Generated draft: requires mathematical and source review.
\documentclass[11pt]{article}
\usepackage[margin=25mm]{geometry}
\usepackage{fontspec}
\setmainfont{Latin Modern Roman}
\usepackage{amsmath,amssymb,mathtools,unicode-math}
\setmathfont{Latin Modern Math}
\usepackage{enumitem,needspace,xurl,hyperref}
\setlength{\parindent}{0pt}
\setlength{\parskip}{5pt}
\setlength{\emergencystretch}{3em}
\newcommand{\N}{\mathbb{N}}
\newcommand{\Z}{\mathbb{Z}}
\newcommand{\Q}{\mathbb{Q}}
\newcommand{\R}{\mathbb{R}}
\newcommand{\C}{\mathbb{C}}
\newcommand{\F}{\mathbb{F}}
\newcommand{\A}{\mathbb{A}}
\newcommand{\PP}{\mathbb{P}}
\newcommand{\E}{\mathbb{E}}
\newcommand{\eps}{\varepsilon}
\newcommand{\dd}{\,\mathrm{d}}
\begin{document}
\section*{UnsolvedMath ID: 20002009}
\textbf{Source problem code:} catalogue entry
\subsection{Definitions and mathematical statement}
Item C.4 contains two recognition problems. Part A asks when a pure \(d\)-complex is a tropical variety. Its ambient space, rationality, weights, coefficient field and whether tropical means balanced or algebraically realizable are not specified. Under an algebraic-realizability reading one would ask whether a given polyhedral support is \(\operatorname{Trop}(V(I))\) for suitable field and ideal; this is a possible target, not a silent resolution of the wording. No answer to Part A is supplied by the paired hyperplane result.

Part B asks for phase images of linear spaces, called ``phlats.'' Repair the damaged torus notation explicitly: for \(L\subset\mathbb C^n\) linear, use \(L^\times=L\cap(\mathbb C^*)^n\) and
\[
 \operatorname{Arg}:(\mathbb C^*)^n\to
 (\mathbb R/2\pi\mathbb Z)^n,\qquad
 z\longmapsto(\arg z_1,\ldots,\arg z_n).
\]
The actual image and its closure are different recognition questions.

The aid treats a single homogeneous hyperplane
\[
 L=\{z:\sum_{j=1}^n c_jz_j=0\},\qquad
 I=\{j:c_j\ne0\}\ne\varnothing.
\]
For a phase \(\theta\), put \(u_j=e^{i(\arg c_j+\theta_j)}\), \(j\in I\). Coordinates outside \(I\) impose no condition. The reported equivalent actual-image tests are
\[
 \theta\in\operatorname{Arg}(L^\times)
 \ \Longleftrightarrow\
 \exists\rho_j>0:\ \sum_{j\in I}\rho_ju_j=0
 \ \Longleftrightarrow\
 0\in\operatorname{relint}\operatorname{conv}\{u_j:j\in I\}.
\]
The linear feasibility test has two real equations, using real and imaginary parts. It may be normalized by \(\sum_j\rho_j=1\).

For closure, the reported tests replace strict positivity by \(\rho_j\geq0\), not all zero, and relative interior by ordinary convex-hull membership:
\[
 \theta\in\overline{\operatorname{Arg}(L^\times)}
 \ \Longleftrightarrow\
 0\in\operatorname{conv}\{u_j:j\in I\}.
\]
The nonzero or normalized requirement avoids the vacuous zero solution.

For the cyclic-gap version, list the distinct effective directions around \(S^1\), including the wrap-around gap, and let \(g_{\max}\) be the maximum gap. Actual membership is equivalent to \(g_{\max}<\pi\), with the sole equality exception of exactly two distinct antipodal directions. Closure membership is equivalent to \(g_{\max}\leq\pi\). Repeated directions are allowed and do not count as additional distinct directions. A sole direction has cyclic gap \(2\pi\). An antipodal pair with an additional distinct direction confined to one closed semicircle is boundary-only, not an actual-image point.

For an affine hyperplane \(\sum c_jz_j=b\), \(b\ne0\), the aid extends the same criterion by adjoining the fixed direction \(-b/|b|\). A positive relation can be rescaled so its constant-direction coefficient is \(|b|\); this is the necessary normalization for the affine equation. The case \(b=0\) is the homogeneous case, not an added zero direction. These are single-hyperplane criteria and do not provide a characterization for all higher-codimension linear spaces or solve tropical recognition in Part A.
\subsection{Short English statement}
The source asks both tropical-complex recognition and phase-image recognition. The supplied aid solves a focused hyperplane phase problem: actual phases require a strictly positive zero relation, while closure allows a nonnegative one. The cyclic-gap equality case must distinguish a pure antipodal pair from boundary-only configurations. The broader two recognition problems remain outside that special case.
\subsection{Sources}
\begin{itemize}
\item Local source record: catalogue entry, catalogue entry; canonical \texttt{aim-geometry-notes.json}, \texttt{notes[346]}, item C.4, ``Recognition problems.'' Title: ``Exact phase recognition for a complex hyperplane.''
\item AIM workshop ``Amoebas and tropical geometry,'' 23--26 October 2003, supplied reference \url{https://aimath.org/WWN/amoebas/amoebas.pdf}. This draft uses the supplied record without accessing a PDF. Original remarks and literature fields are empty.
\item Research attempt 1 explicitly reconstructs the damaged coordinatewise phase map and supplies the homogeneous/affine hyperplane feasibility, convex-hull and cyclic-gap criteria. No authoritative external bibliographic citation is supplied.
\item Registry dated 14 August 2026: ``partially\_solved,'' ``PARTIAL-PROGRESS,'' statement status ``exact''; attempt partial, rigor ``proved,'' novelty confidence low. The machine-generated, non-peer-reviewed aid requires independent verification before citation.
\end{itemize}
\end{document}
